\subsection{The Test-Oracle Problem and LLM-Based Test Generation}
\label{subsec:rw6}

A test exercises a program; an \emph{oracle} decides whether the observed
behaviour is correct. \citet{barr2015oracle} survey this oracle problem and
establish that, while input generation is now largely automatable, the oracle
remains the binding constraint on test automation: without a trustworthy
arbiter of correctness, generated tests can only detect crashes and uncaught
exceptions, not functional deviation. Their taxonomy distinguishes
\emph{specified} oracles, derived from a formal description of intended
behaviour, from \emph{derived} and \emph{implicit} oracles inferred from
artefacts or general expectations. This article positions its formal-contract
layer squarely as a \emph{specified} oracle, but one that is recovered and
ratified rather than assumed to pre-exist.

A long line of work attacks the oracle problem by sidestepping the need for an
absolute correctness criterion. Metamorphic testing checks
\emph{relations} between multiple executions---for example, that permuting a
list and sorting it yields the same result as sorting and permuting---rather
than asserting any single output; \citet{chen2018metamorphic} and
\citet{segura2016metamorphic} survey the technique and its reach into
domains where no conventional oracle exists. Metamorphic relations are
expressive but partial: they constrain behaviour without fully pinning it,
and an implementation can satisfy every relation while still being wrong. A
complementary tradition \emph{infers} oracles from program behaviour. Daikon
detects likely invariants by observing values across executions and reporting
properties that held over the observed runs \citep{ernst2007daikon}; the
resulting invariants are candidate oracles, but they describe what the code
\emph{did}, not what it \emph{should} do, and are unsound by construction. This
description-versus-prescription gap is exactly the tension this article
confronts in its brownfield characterisation phase, where a code-derived
contract is treated as evidence of intent and prime suspect simultaneously.

Search-based and feedback-directed generators sharpened the gap by automating
inputs at scale while leaving the oracle largely implicit. EvoSuite evolves
whole test suites toward coverage targets and, lacking a specification, emits
\emph{regression} assertions that merely capture current behaviour
\citep{fraser2011evosuite}; Randoop generates sequences by feedback-directed
random testing and flags only contract violations and crashes
\citep{pacheco2007randoop}. Both expose the central weakness: a test suite with
self-generated assertions is a snapshot, not a correctness judgment, and it
will faithfully ratify a buggy implementation. Learning-based oracle synthesis
followed, treating assertion generation as a translation task. ATLAS learns to
emit assert statements for a test and its focal method via neural machine
translation \citep{watson2020atlas}, and AthenaTest fine-tunes a transformer on
mined focal-method/test pairs \citep{tufano2020athenatest}. These systems learn
the \emph{form} of plausible oracles from human-written corpora but inherit
whatever (in)correctness those corpora encode, and they offer no machine-
checkable guarantee that a generated assertion reflects intended behaviour.

The arrival of large language models (LLMs) has accelerated this work without
resolving the underlying problem. TestPilot generates and adaptively repairs
unit tests by re-prompting the model with failing tests and error messages
\citep{schafer2024testpilot}; ChatTester iteratively refines ChatGPT-produced
tests and improves over earlier neural baselines on compilability and coverage
\citep{yuan2024chattester}; broad empirical studies report that LLM-generated
suites achieve competitive coverage and superior readability yet still suffer
correctness and compilation defects \citep{siddiq2024empirical}. Crucially,
recent analyses of \emph{LLM-driven oracle generation} document a structural
hazard: when the same model both implements the code and writes the asserting
oracle, the oracle tends to encode the model's own (possibly mistaken)
expectations, yielding tests that pass by construction
\citep{liu2026oracle}. This is the self-grading, circular-verification defect
that motivates the present work.

The gap, then, is not a shortage of test-generation machinery but the absence
of an oracle that is at once \emph{machine-checkable}, \emph{prescriptive}, and
\emph{independently authored}. Sampling-based oracles---metamorphic relations,
inferred invariants, learned or LLM-written assertions---judge correctness on a
finite set of executions and, when produced by the implementing agent,
collapse into circularity. This article instead elevates a formal JML contract
to the oracle and discharges it with OpenJML over \emph{all} inputs rather than
samples, separating the elaborator role that authors the contract from the
implementer role that satisfies it (the independent-authority constraint). Where
prior oracle-inference work stops at unsound, code-derived descriptions, the
proposed pipeline tags each obligation with provenance and assurance,
triangulates code against tickets and documentation to surface bugs as
\emph{disagreements}, and feeds OpenJML counterexamples back into a
CEGIS-shaped synthesis loop---turning the inferred-invariant tradition of
\citet{ernst2007daikon} and the LLM-test tradition of
\citet{schafer2024testpilot} from a description-and-eyeball workflow into a
verified, prescriptive, change-propagating contract layer.
